\chapter{Introducción}
\label{ch:intro}

\section{Motivación}
\label{sec:intro-motivacion}

La generalización de los dispositivos de geolocalización ha transformado el
estudio del movimiento animal \cite{kays2015}. Repositorios como Movebank acumulan
hoy millones de posiciones GPS de miles de individuos, de modo que el cuello de
botella de la
disciplina se ha desplazado de \emph{registrar} dónde ha estado un animal a
\emph{modelar y anticipar} dónde estará. Predecir la posición futura a partir del
rastro histórico es, sin embargo, un problema abierto: el movimiento combina
rutina y cambio, responde a factores ambientales y varía de un individuo a otro.
Esa dificultad motiva combinar modelos probabilísticos y técnicas de aprendizaje
automático en un enfoque híbrido, en lugar de depender de un único tipo de
modelo.

Este trabajo aborda esa pregunta sobre una especie concreta, la gaviota sombría
(\textit{Larus fuscus}), un migrante de larga distancia cuyo corredor se extiende
entre las zonas de cría del norte de Europa y las áreas de invernada del este de
África. Su ciclo anual está marcado por una fenología nítida (cría en junio y
julio, paso migratorio en primavera y otoño e invernada en los meses centrales
del invierno), lo que la convierte en un caso idóneo: una única especie, con un
patrón estacional fuerte y datos de seguimiento GPS públicos
\cite{movebank_gulls, wikelski2015}. El conjunto de datos concreto se describe en
el capítulo \ref{ch:o1}.

\section{Definición del problema y enfoque híbrido}
\label{sec:intro-enfoque}

La pregunta que vertebra el trabajo es deliberadamente simple: dada la posición
de un individuo hoy, ¿dónde estará mañana? Se trata de una predicción a un día
sobre una secuencia de posiciones diarias, la que produce la preparación de los
datos del capítulo \ref{ch:o1}. Plantearla a esa escala (un paso por día) evita
depender de la frecuencia irregular del seguimiento GPS y centra el problema en
el desplazamiento, no en el muestreo.

Predecir ese desplazamiento es difícil porque el movimiento alterna regímenes muy
distintos (días casi inmóviles y días de migración de cientos de kilómetros) y
porque cada ave sigue su propia ruta, de modo que un patrón
aprendido de la población no describe sin más a un individuo concreto. A ello se
añade una tensión metodológica: los modelos probabilísticos clásicos son
interpretables pero limitados, mientras que el aprendizaje automático capta
relaciones no lineales a costa de transparencia. El enfoque híbrido de este
trabajo combina dos piezas para afrontarla: una baseline probabilística
interpretable, frente a la que se mide todo lo demás, y un predictor que une un
modelo de comportamiento con aprendizaje automático para tratar de superarla
(figura \ref{fig:intro-hibrido}).

\begin{figure}[H]
  \centering
  \includegraphics[width=\linewidth]{intro_enfoque_hibrido.pdf}
  \caption[El enfoque híbrido: baseline frente a predictor combinado]{Esquema del
  enfoque híbrido. A partir de la secuencia de posiciones diarias, la cadena de
  Markov visible (vía inferior) actúa como baseline y el predictor híbrido (vía
  superior) combina el modelo oculto de Markov con el aprendizaje supervisado;
  ambas vías se comparan en el visor interactivo de la posición y el error.}
  \label{fig:intro-hibrido}
\end{figure}

La baseline es una cadena de Markov visible, que estima la posición del día
siguiente a partir de la celda actual; es interpretable y fija el criterio de
comparación común al resto de modelos. El predictor híbrido, por su parte,
encadena dos modelos: un modelo oculto de Markov infiere para cada día un estado
de comportamiento no observado (estacionario o en migración), y ese estado alimenta,
como una variable más, a los modelos de aprendizaje supervisado (Random Forest,
XGBoost y LightGBM), que predicen el desplazamiento incorporando además variables
propias de cada individuo. Esa unión de un modelo probabilístico de comportamiento
con el aprendizaje automático es lo que da al enfoque su carácter híbrido y lo que
busca superar a la baseline. Para interpretar los resultados se desarrolla además
una herramienta de cartografía interactiva, que permite visualizar de un vistazo
sobre el mapa las predicciones de los modelos y analizar dónde y cuándo acierta o
falla cada una.

Conviene anticipar el hallazgo que recorre todos los capítulos: la regla trivial
de suponer que mañana el ave seguirá donde está hoy (la persistencia) resulta
sorprendentemente difícil de superar, porque la especie pasa la mayor parte de
los días sin cambiar de lugar. El valor del trabajo no reside, por tanto, en un
salto de exactitud sobre esa regla, sino en calibrar mejor la incertidumbre, en
aislar el régimen de migración (donde la predicción sí aporta) y en caracterizar
el techo estructural de predecir a un día con información local.

\section{Objetivos del trabajo}
\label{sec:intro-objetivos}

El objetivo general del trabajo es diseñar y evaluar un enfoque híbrido, que
combine modelos probabilísticos y aprendizaje supervisado, para predecir el
desplazamiento diario de la gaviota sombría a partir de datos de seguimiento GPS.
Reviste especial interés la comparación sistemática de las distintas soluciones
desarrolladas, medidas todas con una métrica probabilística común y frente a
baselines explícitas.

Este objetivo general se concreta en cinco objetivos específicos:

\begin{enumerate}
  \item Transformar los registros GPS en bruto en una secuencia diaria homogénea
  por individuo, con los huecos del seguimiento representados de forma explícita,
  que sirva de base común a todos los modelos.
  \item Establecer una baseline predictiva interpretable mediante cadenas de
  Markov visibles, y fijar con ella el criterio de evaluación (el log-loss) y la
  referencia de persistencia frente a la que se contrasta el resto de modelos.
  \item Detectar el régimen de comportamiento de cada día (estacionario o en
  migración) mediante un modelo oculto de Markov, comprobando que los estados
  inferidos son coherentes con la fenología de la especie.
  \item Predecir la posición del día siguiente con modelos de aprendizaje
  supervisado (Random Forest, XGBoost y LightGBM) que integran el comportamiento
  inferido y variables propias de cada individuo, empleando solo información
  disponible antes del instante que se predice.
  \item Desarrollar una herramienta de cartografía interactiva que permita
  visualizar y comparar sobre el mapa los resultados de los modelos.
\end{enumerate}
